%%%%%%%%%%%%%%%%%%%%%%%%%%%%%%%%%%%%%%%%%
% baposter Portrait Poster
% LaTeX Template
% Version 1.0 (15/5/13)
%
% Created by:
% Brian Amberg (baposter@brian-amberg.de)
%
% This template has been downloaded from:
% http://www.LaTeXTemplates.com
%
% License:
% CC BY-NC-SA 3.0 (http://creativecommons.org/licenses/by-nc-sa/3.0/)
%
%%%%%%%%%%%%%%%%%%%%%%%%%%%%%%%%%%%%%%%%%

%----------------------------------------------------------------------------------------
%	PACKAGES AND OTHER DOCUMENT CONFIGURATIONS
%----------------------------------------------------------------------------------------

\documentclass[a0paper,portrait]{baposter}

\usepackage[font=small,labelfont=bf]{caption} % Required for specifying captions to tables and figures
\usepackage{booktabs} % Horizontal rules in tables
\usepackage{relsize} % Used for making text smaller in some places

\graphicspath{{figures/}} % Directory in which figures are stored

\definecolor{bordercol}{RGB}{40,40,40} % Border color of content boxes
\definecolor{headercol1}{RGB}{186,215,230} % Background color for the header in the content boxes (left side)
\definecolor{headercol2}{RGB}{80,80,80} % Background color for the header in the content boxes (right side)
\definecolor{headerfontcol}{RGB}{0,0,0} % Text color for the header text in the content boxes
\definecolor{boxcolor}{RGB}{186,215,230} % Background color for the content in the content boxes

\begin{document}

\background{ % Set the background to an image (background.pdf)
\begin{tikzpicture}[remember picture,overlay]
\draw (current page.north west)+(-2em,2em) node[anchor=north west]
{\includegraphics[height=1.1\textheight]{background}};
\end{tikzpicture}
}

\begin{poster}{
grid=false,
borderColor=bordercol, % Border color of content boxes
headerColorOne=headercol1, % Background color for the header in the content boxes (left side)
headerColorTwo=headercol2, % Background color for the header in the content boxes (right side)
headerFontColor=headerfontcol, % Text color for the header text in the content boxes
boxColorOne=boxcolor, % Background color for the content in the content boxes
headershape=roundedright, % Specify the rounded corner in the content box headers
headerfont=\Large\sf\bf, % Font modifiers for the text in the content box headers
textborder=rectangle,
background=user,
headerborder=open, % Change to closed for a line under the content box headers
boxshade=plain
}
{}
%
%----------------------------------------------------------------------------------------
%	TITLE AND AUTHOR NAME
%----------------------------------------------------------------------------------------
%
{\sf\bf Testing the upper-mantle seismic signature of plume- \\ lithosphere interaction in the northern East-African Rift } % Poster title
{\vspace{1em} Chiara Civiero, \textbf{John Armitage}, Sakia Goes, James Hammond\\ % Author names
{\smaller cciviero@cp.dias.ie, armitage@ipgp.fr, saskia.goes@imperial.ac.uk, j.hammond@ucl.ac.uk }} % Author email addresses
{\includegraphics[scale=0.25]{LogoIPGP}} % University/lab logo

%---------------------------------------------------------------------------------------
%	INTRODUCTION
%---------------------------------------------------------------------------------------

\headerbox{Introduction}{name=introduction,column=0,row=0}{

Below the northern East African Rift two low seismic velocity zones are imaged. \textbf{Are these due to the rising thermal anomalies?} We will explore this hypothesis by generating synthetic tomography from models of the convective destabilization of hot mantle.

\begin{center}
\includegraphics[width=0.8\linewidth]{fig01}
\captionof{figure}{The northern East African Rift and slices from the P-wave tomographic model NEAR-P15 \cite{civiero-etal-2015}.}
\end{center}

}

%---------------------------------------------------------------------------------------
%	METHODS
%---------------------------------------------------------------------------------------

\headerbox{Geodynamic to Seismic}{name=methods,column=0,below=introduction}{

\smaller
Both Rayleigh B{\'e}nard convection and Rayleigh Taylor instabilities are generated using by solving for modified Stokes flow in a Cartesian box using the code Stag3D \cite{tackley-1998}. Temperature and pressure are converted to mineral phases assuming a pyrolite composition for the mantle and harzbugite for the lithosphere. This gives the anharmonic seismic velocity. This is corrected for attenuation using the $Q_{g}$ model \cite{goes-etal-2012}.
\begin{center}
\includegraphics[width=0.8\linewidth]{velocities}
\captionof{figure}{2D slices through the model domains for Rayleigh B{\'e}nard convection with a Newtonian rheology, a non-Newtonian rheology, and the Rayleigh Taylor destabilization of a hot layer.}
\end{center}
}

%----------------------------------------------------------------------------------------
%	REFERENCES
%----------------------------------------------------------------------------------------

\headerbox{References}{name=references,column=0,below=methods}{

\smaller % Reduce the font size in this block
\renewcommand{\section}[2]{\vskip 0.05em} % Get rid of the default "References" section title
\nocite{*} % Insert publications even if they are not cited in the poster

\bibliographystyle{abbrv}
\bibliography{sample.bib} % Use sample.bib as the bibliography file
}

%----------------------------------------------------------------------------------------
%	ACKNOWLEDGEMENTS
%----------------------------------------------------------------------------------------

\headerbox{Acknowledgements}{name=acknowledgements,column=0,below=references, above=bottom}{

\smaller % Reduce the font size in this block
This work was funded by a Agence Nationale de la Recherche grant to John Armtiage and a Janet Watson Fellowship awarded to Chiara Civiero.
} 

%----------------------------------------------------------------------------------------
%	RESULTS 1
%----------------------------------------------------------------------------------------

\headerbox{Rayleigh B{\'e}nard or Taylor?}{name=results1,span=2,column=1,row=0}{ % To reduce this block to 1 column width, remove 'span=2'

For Raleigh B{\'e}nard (RB) convection the wavelength between plumes, $\lambda$, varies as $Ra^{1/6}$ \cite{galsa-2007}. For Raleigh Taylor (RT) destabilisation of a hot layer, $\lambda$, varies as a linear function of the model height \cite{schmeling-1987}. The observed spacing of the two low seismic velocity structures is $\sim700$\,km (Figure 1). We find that for RB convection this implies a $10^6 < Ra < 10^7$. For RT destabilisation this is consistent with a model aspect ratio of $2,100\times2,100\times700$\,km.

\begin{center}
\includegraphics[width=0.49\linewidth]{comparison}
\includegraphics[width=0.49\linewidth]{comparison-3x4-lables.png}
\captionof{figure}{RB convection for Newtonian and non-Newtonian creep with a $100\rm\,^{\circ}C$ hot lower boundary condition (left); RT destabilization for a $100\rm\,^{\circ}C$ 100\,km thick hot layer. Aspect ratios of 3x3x1 and 4x4x1 (right)}
\end{center}

}

%----------------------------------------------------------------------------------------
%	RESULTS 2
%----------------------------------------------------------------------------------------

\headerbox{Synthetic Tomography: Rayleigh Taylor wins}{name=results2,span=2,column=1,below=results1}{ % To reduce this block to 1 column width, remove 'span=2'

\begin{enumerate}
\item It was found that while the spacing of Newtonian RB plumes was consistent with observations, the anomalies are too weak to correspond with the magnitudes of the observations.
\item The RT destabilization of a hot layer that is 2,100 wide gives a pattern of seismic anomalies that matches the observations .
\item The RT destabilization of a $200\rm\,^{\circ}C$ hot layer provides a magnitude that is in loser agreement to the observed anomalies.
\end{enumerate}


%------------------------------------------------

\begin{center}
\includegraphics[width=0.49\linewidth]{fig07}
\includegraphics[width=0.49\linewidth]{fig11}
\captionof{figure}{Figure caption 1 (left); Figure caption 2 (right)}
\end{center}

}

\headerbox{Conclusions and Implications}{name=conclusions,span=2,column=1,below=results2,above=bottom}{

\begin{enumerate}
    \item Seismic tomography from the dense array of seismometers in the northern East African Rift is consistent with the Rayleigh Taylor-like destabilization of a hot layer of mantle at or below the 660\,km discontinuity.
    \item The thermal anomalies below Afar and the Main Ethiopian Rift are the form of rising plumes.
    \item This is a strong statement because it means that the plume-like structures are not secondary plumes from a stalled super-plume but the hot material has its origins in the lower mantle. The hot material subsequently destabilized leading to the volcanism we observe in Afar, the Main Ethiopian Rift, and potentially in the Arabian Peninsular.
    \item One possible mechanism for such destabilization of ponded large plumes is internal heating, which can give additional buoyancy to thermo-chemical plumes \cite{fourel-etal-2018}
\end{enumerate}
}
%----------------------------------------------------------------------------------------

\end{poster}

\end{document}